\title{The Era of 1-bit LLMs: All Large Language Models are in 1.58 Bits}

\begin{abstract}
Recent advancements, including BitNet~\cite{bitnet}, are ushering in a new phase of 1-bit Large Language Models (LLMs). In this study, we present a 1-bit LLM variant called \textbf{\text{BitNet b1.58}}, where each individual parameter (or weight) of the LLM takes ternary values \{-1, 0, 1\}. This model matches the performance of full-precision (i.e., FP16 or BF16) Transformer LLMs of equivalent size and training data in terms of perplexity and downstream task accuracy, while offering substantial improvements in latency, memory usage, throughput, and energy efficiency. More importantly, the 1.58-bit LLM establishes a new scaling law and training paradigm for future LLM generations that are both efficient and high-performing. Additionally, it paves the way for a novel computational framework and facilitates the design of specialized hardware tailored for 1-bit LLMs.
\end{abstract}

\section{The Era of 1-bit LLMs}

In recent years, the AI community has witnessed a swift expansion in the scale and abilities of Large Language Models (LLMs). These models have achieved impressive results across diverse natural language processing tasks, yet their growing size introduces deployment difficulties and raises concerns about environmental and financial costs due to elevated energy consumption.

A strategy to tackle these issues is post-training quantization, which creates low-bit models for inference~\cite{smoothquant, gptq, quip, quip_sharp}. This method lowers the precision of weights and activations, substantially cutting down the memory footprint and computational load of LLMs. The trend has shifted from 16-bit precision toward lower bit-widths, such as 4-bit versions~\cite{gptq, awq}. Nevertheless, post-training quantization remains suboptimal, despite its widespread adoption in industrial LLM implementations.

Recent research on 1-bit model architectures, such as BitNet~\cite{bitnet}, highlights a promising approach for lowering the costs associated with LLMs while preserving their effectiveness. Standard LLMs operate using 16-bit floating-point values (e.g., FP16 or BF16), with the bulk of their operations dominated by matrix multiplication. Consequently, the primary computational expense arises from floating-point addition and multiplication. In contrast, BitNet’s matrix multiplication relies solely on integer addition, leading to substantial energy savings for LLMs. Since power consumption often represents the fundamental bottleneck for compute performance in many chips, these energy savings can be converted into faster computation speeds.

Beyond computation, transferring model parameters from DRAM to the on-chip accelerator memory (such as SRAM) during inference can be costly. Efforts to expand SRAM capacity to boost throughput have been made but come at a significantly higher cost compared to DRAM. Relative to full-precision models, 1-bit LLMs require considerably less memory in terms of both capacity and bandwidth, which can drastically decrease the expense and duration of loading weights from DRAM, resulting in quicker and more efficient inference.

In this study, we propose a notable 1-bit LLM variant named \textbf{\text{BitNet b1.58}}, where every parameter is ternary, adopting values from \{-1, 0, 1\}. By introducing the value 0 to the original 1-bit BitNet scheme, we achieve an effective precision of 1.58 bits in binary representation. \text{BitNet b1.58} preserves all advantages of the original 1-bit BitNet, including its innovative computation method that almost entirely eliminates multiplication in matrix multiplication and is highly amenable to optimization. Moreover, it matches the original 1-bit BitNet in energy consumption while being substantially more efficient in memory usage, throughput, and latency when compared to FP16 LLM baselines. Additionally, \text{BitNet b1.58} offers two further benefits. First, its modeling power is enhanced due to explicit feature filtering enabled by the inclusion of 0 in model weights, which can significantly elevate the performance of 1-bit LLMs. Second, our experimental results indicate that \text{BitNet b1.58} can attain parity with full-precision (i.e., FP16) baselines regarding both perplexity and downstream task performance, starting from models of size 3B, when using identical configurations (e.g., model size, training tokens, etc.).

\section{BitNet b1.58}

\text{BitNet b1.58} is built upon the BitNet framework, a Transformer model that substitutes \emph{nn.Linear} layers with \emph{BitLinear}. It is trained from the ground up, utilizing 1.58-bit weights along with 8-bit activations. Compared to the original BitNet, several alterations have been incorporated, which we outline below.

\paragraph{Quantization Function.} To restrict the weights to values of -1, 0, or +1, we employ an \emph{absmean} quantization approach. This method first normalizes the weight matrix by its mean absolute value, and then rounds each element to the nearest integer within \{-1, 0, +1\}:

\begin{equation}
    \widetilde{W} = \mathrm{RoundClip}\left(\frac{W}{\gamma + \epsilon}, -1, 1\right),
\end{equation}

\begin{equation}
    \mathrm{RoundClip}(x, a, b) = \max(a, \min(b, \mathrm{round}(x))),
\end{equation}

\begin{equation}
    \gamma = \frac{1}{nm}\sum_{ij} |W_{ij}|.
\end{equation}

The activation quantization follows the same procedure as in BitNet, except that we omit scaling the activations before non-linear operations to the range $[0, Q_b]$. Instead, each token’s activations are scaled to $[-Q_b, Q_b]$, eliminating the need for zero-point quantization. This approach simplifies implementation and system-level optimization, while having negligible impact on performance as confirmed by our experiments.

\paragraph{LLaMA-alike Components.}

LLaMA~\cite{llama, llama2} has become the standard backbone architecture for many open-source large language models. To align with the open-source community, \text{BitNet b1.58} integrates LLaMA-inspired components. In particular, it employs RMSNorm~\cite{rmsnorm}, SwiGLU~\cite{swiglu}, rotary embeddings~\cite{rope}, and excludes all biases. This design enables easy integration of \text{BitNet b1.58} into widely-used open-source platforms such as Huggingface, vLLM~\cite{vllm}, and llama.cpp\footnote{\url{https://github.com/ggerganov/llama.cpp}} with minimal adaptation.

\section{Results}

We conducted a comparison between \text{BitNet b1.58} and our reproduced FP16 \text{LLaMA LLM} across multiple model sizes. To maintain fairness, both models were pre-trained on the RedPajama dataset~\cite{redpajama} using 100 billion tokens. We assessed their zero-shot capabilities on various language benchmarks, including ARC-Easy~\cite{arc}, ARC-Challenge~\cite{arc}, Hellaswag~\cite{hellaswag}, Winogrande~\cite{winoGrande}, PIQA~\cite{piqa}, OpenbookQA~\cite{openbookqa}, and BoolQ~\cite{boolq}. Additionally, validation perplexity was measured on the WikiText2~\cite{wikitext2} and C4~\cite{c4} datasets.

We also compared the GPU memory usage and latency during runtime for both \text{LLaMA LLM} and \text{BitNet b1.58}. These metrics were obtained using the FasterTransformer\footnote{\url{https://github.com/NVIDIA/FasterTransformer}} framework, which is highly optimized for LLM inference on GPUs. The 2-bit kernel from Ladder~\cite{ladder} was incorporated for \text{BitNet b1.58}. We reported inference time per output token, as it constitutes the primary cost during inference.

Table~\ref{tab:ppl} provides a summary of the perplexity scores and computational costs for \text{BitNet b1.58} and \text{LLaMA LLM}. The table indicates that \text{BitNet b1.58} begins to match the full-precision \text{LLaMA LLM} at the 3B parameter scale in terms of perplexity, while being 2.71 times faster and requiring 3.55 times less GPU memory. Notably, the 3.9B variant of \text{BitNet b1.58} is 2.4 times faster and uses 3.32 times less memory, yet substantially outperforms the \text{LLaMA LLM} 3B model.

Detailed zero-shot accuracy results on downstream tasks are presented in Table~\ref{tab:zero-shot}. We employed the evaluation pipeline from \emph{lm-evaluation-harness}\footnote{\url{https://github.com/EleutherAI/lm-evaluation-harness}}. The findings illustrate that the performance gap between \text{BitNet b1.58} and \text{LLaMA LLM} decreases as model size grows. Crucially, starting from the 3B scale, \text{BitNet b1.58} matches the performance of the full-precision baseline. Consistent with the perplexity observations, the end-task results confirm that \text{BitNet b1.58} 3.9B surpasses \text{LLaMA LLM} 3B while incurring lower memory and latency costs. This establishes \text{BitNet b1.58} as a Pareto improvement relative to the leading LLM models.

\paragraph{Memory and Latency}

We extended the model scaling to sizes 7B, 13B, and 70B and assessed the associated costs. Figure~\ref{fig:latency-memory-energy} presents the patterns of latency and memory usage, demonstrating that the acceleration improves as the model size increases. Notably, \text{BitNet b1.58} 70B achieves a speed-up that is 4.1 times greater than the \text{LLaMA LLM} baseline. This improvement arises because the time required for \emph{nn.Linear} operations increases with model size. Memory usage exhibits a similar behavior, since the embedding remains in full precision and its share of memory diminishes as models become larger. Both latency and memory measurements were conducted using a 2-bit kernel, indicating potential for further cost reductions through optimization.

\paragraph{Energy}

We additionally evaluated the energy consumption related to arithmetic operations for both \text{BitNet b1.58} and \text{LLaMA LLM}. Our primary focus was on matrix multiplication calculations, as they constitute the largest part of LLM computational costs. Figure~\ref{fig:energy} depicts the breakdown of energy expenditure. Most of the energy consumption in \text{BitNet b1.58} is due to INT8 addition operations, whereas \text{LLaMA LLM} requires both FP16 addition and FP16 multiplication. Based on the energy model from~\cite{energycost, pokebnn}, \text{BitNet b1.58} reduces the arithmetic operations energy for matrix multiplication by a factor of 71.4 on 7nm chips. We also presented the total energy usage for processing 512 tokens. Our findings indicate that as model size grows, \text{BitNet b1.58} becomes increasingly energy efficient relative to the FP16 \text{LLaMA LLM} baseline. This efficiency gain is attributed to the increasing proportion of \emph{nn.Linear} operations with model size, while the energy costs of other components become less significant in larger models.

\paragraph{Throughput}

We evaluate the throughput of \text{BitNet b1.58} and \text{LLaMA LLM} with 70 billion parameters on a setup of two 80GB A100 GPUs, utilizing pipeline parallelism~\cite{gpipe} to enable running \text{LLaMA LLM} 70B on these devices. The batch size was progressively increased until the GPU memory capacity was maxed out, using a sequence length of 512. As shown in Table~\ref{tab:throughput}, \text{BitNet b1.58} 70B can handle batch sizes up to 11 times larger than \text{LLaMA LLM}, leading to an 8.9-fold improvement in throughput.

\textbf{\text{BitNet b1.58} establishes a new scaling relationship between model performance and inference cost}. For reference, based on the outcomes presented in Figures~\ref{fig:latency-memory-energy} and \ref{fig:energy}, the equivalences between various model sizes in 1.58-bit and 16-bit formats are as follows:

\begin{itemize}
    \item A 13B \text{BitNet b1.58} model outperforms a 3B FP16 LLM in terms of latency, memory consumption, and energy efficiency.
    \item A 30B \text{BitNet b1.58} model is more efficient than a 7B FP16 LLM with regard to latency, memory usage, and energy usage.
    \item A 70B \text{BitNet b1.58} model surpasses a 13B FP16 LLM in latency, memory, and energy performance.
\end{itemize}

\paragraph{Training with 2T Tokens}

The quantity of training tokens is a key factor for LLMs. To evaluate the scalability of \text{BitNet b1.58} with respect to token count, we trained a \text{BitNet b1.58} model on 2 trillion tokens following the data preparation method from StableLM-3B~\cite{StableLM-3B-4E1T}, which represents the current state-of-the-art open-source 3B model. Both models were assessed using a benchmark composed of Winogrande~\cite{winoGrande}, PIQA~\cite{piqa}, SciQ~\cite{sciq}, LAMBADA~\cite{lambada}, and ARC-easy~\cite{arc}. The zero-shot accuracy results are presented in Table~\ref{tab:2t}. For tasks measured by accuracy and normalized accuracy, we computed the average of the two. The StableLM 3B results at 2T tokens are directly sourced from its technical report. Our results demonstrate that \text{BitNet b1.58} attains superior performance across all evaluation tasks, suggesting that 1.58-bit LLMs possess strong generalization abilities as well.

\section{Discussion and Future Work}

\textbf{1-bit Mixture-of-Experts (MoE) LLMs}

Mixture-of-Experts (MoE) models have shown to be an economical solution for large language models (LLMs). Although they substantially decrease computational FLOPs, their high memory demands and inter-chip communication costs restrict their deployment and practical use. These issues can be mitigated through 1.58-bit LLMs. Primarily, the smaller memory footprint decreases the number of devices necessary for deploying MoE architectures. Additionally, it greatly lowers the overhead associated with transferring activations over networks. Ultimately, if the entire model fits on a single chip, this overhead would be eliminated.

\textbf{Native Support of Long Sequence in LLMs}

Handling long sequences has emerged as a crucial capability in the era of LLMs. A key obstacle for long-sequence inference lies in the memory consumed by the KV caches. \text{BitNet b1.58} marks a significant advancement towards native long sequence support, as it cuts down activation bit-width from 16 to 8 bits, enabling the doubling of context length under the same resource constraints. This can be further losslessly compressed to 4 bits or even lower in 1.58-bit LLMs, which remains an avenue for future exploration.

\textbf{LLMs on Edge and Mobile}

Deploying 1.58-bit LLMs holds considerable promise for enhancing language model performance on edge and mobile devices. These devices typically face limitations in memory capacity and computational power, which constrain the scale and effectiveness of LLMs. However, the lowered memory and energy demands of 1.58-bit LLMs enable their operation on such platforms, unlocking a variety of applications that were previously unfeasible. This advancement can significantly boost the functionality of edge and mobile devices and foster novel uses of LLMs. Furthermore, 1.58-bit LLMs are more compatible with CPUs, which are the predominant processors in edge and mobile hardware. Thus, \text{BitNet b1.58} can be run efficiently on these devices, further enhancing their performance and capabilities.

\textbf{New Hardware for 1-bit LLMs}

Recent developments like~Groq\footnote{\url{https://groq.com/}} have shown encouraging outcomes and potential for developing specialized hardware (e.g., LPUs) tailored for LLMs. Moving forward, we foresee and advocate for the creation of novel hardware and system designs specifically optimized for 1-bit LLMs, enabled by the novel computational paradigm introduced in BitNet~\cite{bitnet}.

\end{document}